\documentclass[11pt]{amsart}
\usepackage[T1]{fontenc}
\usepackage[utf8]{inputenc}
\usepackage{lmodern}
\usepackage{amsmath,amssymb,amsthm,mathtools}
\usepackage{microtype}
\usepackage{enumitem}
\usepackage[hidelinks]{hyperref}
\usepackage[a4paper,margin=32mm]{geometry}

\newtheorem{theorem}{Theorem}[section]
\newtheorem{lemma}[theorem]{Lemma}
\newtheorem{proposition}[theorem]{Proposition}
\newtheorem{corollary}[theorem]{Corollary}
\theoremstyle{definition}
\newtheorem{definition}[theorem]{Definition}
\newtheorem{conjecture}[theorem]{Conjecture}
\newtheorem{claim}[theorem]{Claim}

\newcommand{\Z}{\mathbb Z}
\newcommand{\Q}{\mathbb Q}
\newcommand{\R}{\mathbb R}
\newcommand{\F}{\mathbb F}
\newcommand{\Aut}{\operatorname{Aut}}
\newcommand{\Inn}{\operatorname{Inn}}
\newcommand{\Out}{\operatorname{Out}}
\newcommand{\Tor}{\operatorname{Tor}}
\newcommand{\rank}{\operatorname{rank}}
\newcommand{\core}{\operatorname{core}}

\title{An irreducible elementary net over an algebraic extension that is not closed}
\author{GroupTheory50}
\date{August 27, 2026}
\subjclass[2020]{20H25, 20G35}
\keywords{elementary net, carpet subgroup, transvection, principal ideal domain, Kourovka Notebook}

\begin{document}
\begin{abstract}
For every order $n\geq3$ we construct an irreducible elementary net of modules over a principal ideal domain in an algebraic extension of its fraction field that is not closed.  The example is defined over $\mathbb Z$ inside the cubic field generated by a root of $x^3-x-1$.  This gives a negative answer to Kourovka Problem 19.48.
\end{abstract}
\maketitle

\section{Introduction}
Elementary nets, also called carpets, organize the allowed parameters of elementary transvections in matrix groups.  Kourovka Problem 19.48, proposed by V. A. Koibaev, asks whether an irreducible elementary net of $R$-modules over an algebraic extension of the fraction field of a principal ideal domain must be closed \cite{kourovka}.  We give an explicit counterexample.

For pairwise distinct indices $i,r,j$, an elementary net $\sigma=(\sigma_{ij})$ satisfies
\[
 \sigma_{ir}\sigma_{rj}\subseteq\sigma_{ij}.
\]
It is irreducible if every off-diagonal entry is nonzero.  It is closed if the elementary subgroup $E(\sigma)$ contains no transvection $x_{ij}(a)$ with $a\notin\sigma_{ij}$.

\section{The cubic example}
Let $\alpha$ be a root of
\[
 \alpha^3-\alpha-1=0,
\]
put $K=\Q(\alpha)$ and $T=\Z[\alpha]$.  Since
\[
 \alpha^{-1}=\alpha^2-1,
\]
the element $\alpha$ is a unit in $T$.  Define additive subgroups
\[
 A=\Z\oplus\Z\alpha\oplus2\Z\alpha^2,
\]
\[
 B=\Z\oplus2\Z\alpha\oplus\Z(\alpha^2-1),
 \qquad C=2T.
\]
Then $C\subseteq A\cap B$ and
\[
 AC\subseteq C,\qquad BC\subseteq C,\qquad C^2\subseteq C.
\]

For $n\geq3$, set
\[
 \sigma_{12}=A,\qquad \sigma_{21}=B,
\]
and let every other off-diagonal entry be $C$.

\begin{lemma}\label{lem:elementary-net}
The family $\sigma$ is an irreducible elementary net.
\end{lemma}
\begin{proof}
All entries are nonzero.  Along a path $i\to r\to j$ with $i,r,j$ pairwise distinct, the two exceptional entries $\sigma_{12}$ and $\sigma_{21}$ cannot occur consecutively.  Thus every product required by the net condition is one of $AC$, $BC$, or $C^2$, or is contained in an exceptional target containing $C$.  The displayed inclusions give the claim.
\end{proof}

\section{Failure of closure}
In the $(1,2)$ block write
\[
 x(t)=x_{12}(t),\qquad y(s)=x_{21}(s),
\]
and for a unit $t$ put
\[
 w(t)=x(t)y(-t^{-1})x(t)
 =\begin{pmatrix}0&t\\-t^{-1}&0\end{pmatrix}.
\]
We have $\alpha,-1\in A$ and $-\alpha^{-1},1\in B$.  Hence
\[
 w(\alpha)w(-1)\in E(\sigma).
\]
A direct multiplication gives
\[
 w(\alpha)w(-1)=\operatorname{diag}(\alpha,\alpha^{-1})
\]
in the $(1,2)$ block.  Conjugating $x_{12}(1)$ by this diagonal matrix yields
\[
 x_{12}(\alpha^2)\in E(\sigma).
\]
However $\alpha^2\notin A$, because the coefficient of $\alpha^2$ in an element of $A$ is even.

\begin{theorem}\label{thm:irreducible-not-closed}
For every $n\geq3$, the elementary net $\sigma$ above is irreducible but not closed.
\end{theorem}
\begin{proof}
Irreducibility was established in Lemma~\ref{lem:elementary-net}, while $x_{12}(\alpha^2)\in E(\sigma)$ with $\alpha^2\notin\sigma_{12}$ is exactly the failure of closure.
\end{proof}

\begin{corollary}\label{cor:kourovka-19-48}
By Theorem~\ref{thm:irreducible-not-closed}, Kourovka Problem 19.48 has a negative answer, already for the PID $R=\Z$ and the cubic algebraic extension $K=\Q(\alpha)$.
\end{corollary}

\section{Perspective}
The mechanism is a rank-two diagonal element manufactured from allowed elementary transvections.  Its conjugation rescales an allowed parameter outside the prescribed module.  This provides a compact obstruction to closure that persists in every rank $n\geq3$.

\begin{thebibliography}{99}
\bibitem{kourovka} E. I. Khukhro and V. D. Mazurov (eds.), \emph{Unsolved Problems in Group Theory: The Kourovka Notebook}, No. 21, Novosibirsk, 2026; arXiv:1401.0300.
\bibitem{steinberg} R. Steinberg, \emph{Lectures on Chevalley Groups}, Yale University, 1968.
\end{thebibliography}
\end{document}
